\clearpage

\Needspace{12\baselineskip}
\section{API 阅读方法与覆盖边界}
\lead{AC-01 · Core 接口契约}

\subsection{阅读路径与例子的前提}
先按前 35 章找到组件，再按 AC 编号查本节契约。每个签名上方给出所属符号，避免把两个类同名的 cancel 当成同一操作。下方的参数、状态和例子解释如何使用。例子省略证书、Face、配置和业务序列化的建立，不是可直接启动的完整程序；具体类型和重载用 API ID 到 api/ 声明参考查找。

\subsection{常用术语}

\begin{longtable}{@{}p{43mm}p{119mm}@{}}
\toprule
项目 & 行为与调用含义 \\
\midrule\endhead
owner & 负责创建、修改和销毁某项状态的组件；不是仅持有一个指针的调用方。 \\
ACK\_CLOSED & ACK 收集结束后冻结的候选快照。ackClosedDigest 将规划输入绑定到这一次观察；不能换成模型摘要。 \\
proposal / plan core & proposal 是待验证的放置建议；plan core 是封存的执行内容。产生建议不等于获准执行。 \\
projection / grant view & projection 是给特定参与者的选择信息；grant view 是授权计算需要的计划视图；两者不是任意可互换 JSON。 \\
recipe / exact reuse & recipe 描述怎样装配实际制品。exact reuse 要求身份、角色、设备和有效期均匹配，不能只看 has\_model。 \\
epoch / lineage & epoch 必须带所属域理解；lineage 关联生成前缀与执行轮次。Controller、Repo、attempt、stream、conversation 的计数不能直接比较。 \\
观察状态 & SIGNATURE\_ONLY 仅有声明；CONTRACT\_REFERENCED 有人工说明引用，仍不等于全部运行行为已验收。 \\
\bottomrule
\end{longtable}

\subsection{一次查阅示例}
需要停止生成时，先查 AC-14 的请求句柄 cancel，再查 AC-18 的模型状态机和 Provider KV 操作。句柄终态取消可无操作，而 Qwen 状态机终态再次 cancel 会抛异常；KV release 是另一项显式资源动作。这样的差异必须由所属符号和状态契约判断，不能由动词猜测。

完整声明查询：\path{api/target-core-reference.md}。

\Needspace{12\baselineskip}
\section{Core API：容器与角色生命周期}
\lead{AC-02 · Core 接口契约}

\subsection{创建与所有权}
ServiceContainer 组合进程中的角色。RuntimeConfig 的 containerIdentity 是容器身份，groupPrefix 是组命名范围，controllerPrefix 指向权限控制器，trustSchemaPath 指定信任规则。\code{addController/addUser/addProvider} 登记角色；use 系列使用外部已有对象，外部对象必须比容器的相关使用期更长。本地 addLocalService 的 RequestT/ResponseT 是 C++ 类型约定，不产生网络服务权限。

\subsection{调用次序}

\begin{Verbatim}[fontsize=\footnotesize,breaklines,breakanywhere]
// 示意：roles/config/handlers 已由应用建立
// 1. 建立身份、Face 与配置
// 2. 登记 Controller/User/Provider 与本地服务
// 3. 登记启动/停止 hook
container.start();
// 4. 应用运行 Face 事件循环，处理服务回调
container.stop();
// 5. 排空应用 worker 后释放借用的角色、Face 和回调资源
\end{Verbatim}

\subsection{前置条件与失败}
受保护的注册方法调用 ensureNotStarted，启动后修改会被拒绝；空角色对象和缺失角色是本地配置错误，不是网络超时。start/stop 的 hook 不能代替应用 worker 的排空。容器本地调用不经过网络权限握手，只有已经可信的进程内调用方才可使用。扩展应把模型加载、数据库连接和飞控 worker 放在应用 hook 中，Core 容器不解释这些领域状态。

\subsection{完整签名与源码定位}

\Needspace{7\baselineskip}
\textbf{API-ee68261266a4}\quad public\par
\textbf{所属符号：}\code{ndn_service_framework::ServiceContainer::addLocalService}\par
\begin{Verbatim}[fontsize=\footnotesize,breaklines,breakanywhere]
template <typename RequestT, typename ResponseT, typename HandlerT>
void
  addLocalService(const ndn::Name& serviceName, HandlerT&& handler)
\end{Verbatim}
\textbf{源码：}\path{ndn-service-framework/ServiceContainer.hpp}，第 73 行。\par

\Needspace{7\baselineskip}
\textbf{API-04ad67003514}\quad public\par
\textbf{所属符号：}\code{ndn_service_framework::ServiceContainer::start}\par
\begin{Verbatim}[fontsize=\footnotesize,breaklines,breakanywhere]
void
  start()
\end{Verbatim}
\textbf{源码：}\path{ndn-service-framework/ServiceContainer.hpp}，第 331 行。\par

\Needspace{7\baselineskip}
\textbf{API-cd596b18ee08}\quad public\par
\textbf{所属符号：}\code{ndn_service_framework::ServiceContainer::stop}\par
\begin{Verbatim}[fontsize=\footnotesize,breaklines,breakanywhere]
void
  stop()
\end{Verbatim}
\textbf{源码：}\path{ndn-service-framework/ServiceContainer.hpp}，第 365 行。\par

完整声明查询：\path{api/target-core-reference.md}。

\Needspace{12\baselineskip}
\section{Core API：请求发起与 Targeted 调用}
\lead{AC-03 · Core 接口契约}

\subsection{请求标识与参数}
普通 RequestService 的 serviceName 表示服务，requestMessage 是序列化请求，timeoutMs 是毫秒预算，strategy 默认 FirstResponding。另一组重载额外提供 ACK 窗口与候选回调。RequestServiceTargeted 额外指定 provider。返回 ndn::Name 是请求标识，后续结果经 ResponseHandler 交付；TimeoutHandler 处理超时，不是返回值的一部分。保持回调捕获资源有效，并按 User 的 Face/I/O 执行上下文调用。

\subsection{最小调用示意}

\begin{Verbatim}[fontsize=\footnotesize,breaklines,breakanywhere]
// user、RequestMessage payload、onTimeout/onResponse 已建立
auto requestId = user.RequestService(
    ndn::Name("/example/echo"), payload, 3000, onTimeout, onResponse);
// 保存 requestId 用于观测；此处尚未取得业务响应
// 已知 provider 的调用使用 RequestServiceTargeted(provider, ...)
\end{Verbatim}

\subsection{事件与接受边界}

\begin{longtable}{@{}p{43mm}p{119mm}@{}}
\toprule
项目 & 行为与调用含义 \\
\midrule\endhead
Request → ACK & Provider 验证发现请求并报价，ACK 说明候选愿意参与；一般 ACK 不保证已独占预留资源。 \\
ACK → Selection & User 按候选策略选择，绑定请求与 token；未被选中的 Provider 不执行请求。 \\
Selection → Response & Provider 验证选择后执行；User 验证响应后才交付。拿到 requestId 或收到 ACK 都不能记为完成。 \\
Targeted & 已有认证 token 关联时走短路径；token 不可用时遵循 bootstrap/失败分支，不跳过授权。 \\
错误观察 & 无 ACK、权限/响应验证失败、业务错误与 deadline 是不同边界。getRequestStatus 是本地状态查询，不能替代远端实际执行证据。 \\
\bottomrule
\end{longtable}

\subsection{回调类型}
低层响应回调收到 ResponseMessage 的常量引用，超时回调收到请求 Name；要异步保存结果就复制所需数据，不捕获一个仅在回调期间有效的引用。类型化重载另行解码 ResponseT。
\begin{Verbatim}[fontsize=\footnotesize,breaklines,breakanywhere]
using ResponseHandler = std::function<void(const ResponseMessage&)>;
using TimeoutHandler = std::function<void(const ndn::Name&)>;
\end{Verbatim}

\subsection{完整签名与源码定位}

\Needspace{7\baselineskip}
\textbf{API-09fc656cb649}\quad public\par
\textbf{所属符号：}\code{ndn_service_framework::ServiceUser::RequestService}\par
\begin{Verbatim}[fontsize=\footnotesize,breaklines,breakanywhere]
ndn::Name RequestService(const ndn::Name& serviceName,
                                 ndn_service_framework::RequestMessage requestMessage,
                                 int timeoutMs,
                                 TimeoutHandler onTimeout,
                                 ResponseHandler onResponseHandler,
                                 size_t strategy = ndn_service_framework::tlv::FirstResponding);
\end{Verbatim}
\textbf{源码：}\path{ndn-service-framework/ServiceUser.hpp}，第 725 行。\par

\Needspace{7\baselineskip}
\textbf{API-1009e77682c2}\quad public\par
\textbf{所属符号：}\code{ndn_service_framework::ServiceUser::RequestServiceTargeted}\par
\begin{Verbatim}[fontsize=\footnotesize,breaklines,breakanywhere]
ndn::Name RequestServiceTargeted(const ndn::Name& provider,
                                 const ndn::Name& serviceName,
                                 ndn_service_framework::RequestMessage requestMessage,
                                 int timeoutMs,
                                 TimeoutHandler onTimeout,
                                 ResponseHandler onResponseHandler);
\end{Verbatim}
\textbf{源码：}\path{ndn-service-framework/ServiceUser.hpp}，第 718 行。\par

\Needspace{7\baselineskip}
\textbf{API-efa9e29aae17}\quad public\par
\textbf{所属符号：}\code{ndn_service_framework::ServiceUser::RequestServiceTargeted}\par
\begin{Verbatim}[fontsize=\footnotesize,breaklines,breakanywhere]
template<typename RequestT, typename ResponseT>
ndn::Name RequestServiceTargeted(const ndn::Name& provider,
                                           const ndn::Name& serviceName,
                                           const RequestT& request,
                                           std::function<void(const ResponseT&)> onResponse,
                                           std::function<void()> onTimeout,
                                           int timeoutMs)
\end{Verbatim}
\textbf{源码：}\path{ndn-service-framework/ServiceUser.hpp}，第 977 行。\par

完整声明查询：\path{api/target-core-reference.md}。

\Needspace{12\baselineskip}
\section{Core API：Provider 注册与作用域回收}
\lead{AC-04 · Core 接口契约}

\subsection{注册者必须保存句柄}
addScopedService 的 ackHandler 处理发现阶段资格，requestHandler 处理被选定后的请求；\code{addScopedCollaborationHandler} 的 handler 接收协作角色与 assignment。返回的 ServiceRegistration 是 move-only 句柄，不能当临时返回值立即丢弃，否则作用域回收会关闭这一代注册。模板服务输入/输出必须满足对应序列化契约。

\subsection{调用与关闭示意}

\begin{Verbatim}[fontsize=\footnotesize,breaklines,breakanywhere]
// service、ackHandler、requestHandler 和模式参数采用上方重载
auto registration = provider.addScopedService(
    service, ackHandler, requestHandler /* 其余参数见签名 */);
// 保存 registration，运行事件循环
registration.close();
// Provider generation 阻止这一代新的 ACK/Selection/执行工作
\end{Verbatim}

\subsection{时序与资源}
在 Face 线程或事件循环开始前注册。Request 认证通过才触发 ACK，Selection 绑定通过才触发执行；不要在 ACK 回调里提前执行有副作用的业务。close/析构先关闭 generation，再投递清理；Provider 已销毁时旧句柄关闭为无害操作。正在运行的外部业务 worker 仍须自行停机/排空，关闭服务不会删除网络缓存中的历史 Data。legacy 注册不能覆盖仍存活的 scoped owner。

\subsection{完整签名与源码定位}

\Needspace{7\baselineskip}
\textbf{API-c3a3fd9da1b0}\quad public\par
\textbf{所属符号：}\code{ndn_service_framework::ServiceProvider::addScopedService}\par
\begin{Verbatim}[fontsize=\footnotesize,breaklines,breakanywhere]
ServiceRegistration addScopedService(const ndn::Name& serviceName,
                                                 AckStrategyHandler ackHandler,
                                                 RequestHandler requestHandler,
                                                 ServiceInvocationMode invocationMode);
\end{Verbatim}
\textbf{源码：}\path{ndn-service-framework/ServiceProvider.hpp}，第 856 行。\par

\Needspace{7\baselineskip}
\textbf{API-23e6ea05d4d6}\quad public\par
\textbf{所属符号：}\code{ndn_service_framework::ServiceProvider::addScopedService}\par
\begin{Verbatim}[fontsize=\footnotesize,breaklines,breakanywhere]
ServiceRegistration addScopedService(const ndn::Name& serviceName,
                                                 AckStrategyHandler ackHandler,
                                                 RequestHandler requestHandler,
                                                 ServiceMode mode);
\end{Verbatim}
\textbf{源码：}\path{ndn-service-framework/ServiceProvider.hpp}，第 860 行。\par

\Needspace{7\baselineskip}
\textbf{API-33314e8f5bbb}\quad public\par
\textbf{所属符号：}\code{ndn_service_framework::ServiceProvider::addScopedCollaborationHandler}\par
\begin{Verbatim}[fontsize=\footnotesize,breaklines,breakanywhere]
ServiceRegistration addScopedCollaborationHandler(
                const ndn::Name& serviceName,
                std::vector<CollaborationRole> allowedRoles,
                AckStrategyHandler ackHandler,
                CollaborationHandler handler);
\end{Verbatim}
\textbf{源码：}\path{ndn-service-framework/ServiceProvider.hpp}，第 869 行。\par

完整声明查询：\path{api/target-core-reference.md}。

\Needspace{12\baselineskip}
\section{Core API：延迟规划与协作提交}
\lead{AC-05 · Core 接口契约}

\subsection{两种时间预算}
ackCollectionTimeMs 只控制报价收集窗口，timeoutMs 控制整个请求期限。例如总预算 5000 ms、ACK 窗口 1000 ms，报价结束后准备与执行共用剩余预算，不能再获得全新的 5000 ms。onAckClosed 给出不可变候选闭包，应用此时才决定本次角色和参与者。

\subsection{两角色例子}
设 role A 在 Provider P1 产生中间对象，role B 在 P2 消费该对象并返回终端结果。应用计划需表达两个角色的参与者、A→B 的依赖及其命名/密钥作用域、终端接收规则；Core 认证和传递这些内容，DI 决定张量是否合法，UAV 决定任务分区是否合法。两个角色不是两个独立、互不绑定的 RequestService。

\subsection{提交顺序与结果}

\begin{Verbatim}[fontsize=\footnotesize,breaklines,breakanywhere]
// 回调顺序示意，不展开应用计划编码
BeginCollaboration(..., onAckClosed, onFinalResponse, onTimeout);
// onAckClosed(snapshot):
//   plan = 按 snapshot 校验候选并构造应用计划
//   accepted = CommitCollaborationPlan(requestId, snapshotDigest, plan)
// accepted 只表示本地提交被接受；等待最终结果或失败
\end{Verbatim}

\subsection{不可越过的状态}
错误 requestId、闭包摘要不匹配、请求过期或状态不允许提交时，不得把失败提交继续显示为执行中。bool 不是各 Provider 的完成聚合值。\code{GenericSelectionTxnStore} 的 prepare/commit 保存 Provider 侧认证选择和 opaque blob，与 User 的 CommitCollaborationPlan 是两个 owner；恢复事务还需匹配 Provider boot epoch、token/lease 和期限，不能用旧 WAL 绕过新进程验证。

\subsection{CollaborationPlan 的具体字段}

\begin{longtable}{@{}p{43mm}p{119mm}@{}}
\toprule
项目 & 行为与调用含义 \\
\midrule\endhead
窗口/期限 & ackCollectionTimeMs 默认 200，timeoutMs 默认 5000，单位毫秒；与 BeginCollaboration 的请求预算一致管理。 \\
roles & 每项含 \code{role/service/requiredArtifact}、\code{allowDynamicProvisioning}=false、provisioningTimeoutMs=30000、\code{appRequirement/assignmentPayload}、\code{minProviders/maxProviders} 默认 1、terminalResponseOwner=false。流式协作指定唯一用户终端 owner。 \\
keyScopes & 每项含 name 与允许的 roles；依赖引用该密钥作用域，而不是自行传递一个未绑定的解密密钥。 \\
dependencies & 每项含 producers、consumers、keyScope、topicPrefix、required=true；A→B 的例子分别在 producers/consumers 填 A/B。 \\
\code{sharedAssignmentMetadata} & 选中参与者共享的不透明 metadata，与各角色 assignmentPayload 分开传递。 \\
participantSelector & shared\_ptr<const \code{ParticipantSelectionPolicy}>；select(candidates, roles) 返回 SelectedParticipant，含精确 Provider、角色、assignment 与实际 ACK 绑定。 \\
\bottomrule
\end{longtable}

\subsection{完整签名与源码定位}

\Needspace{7\baselineskip}
\textbf{API-c9045b5b03e6}\quad public\par
\textbf{所属符号：}\code{ndn_service_framework::ServiceUser::BeginCollaboration}\par
\begin{Verbatim}[fontsize=\footnotesize,breaklines,breakanywhere]
ndn::Name BeginCollaboration(const ServiceName& service,
                                         const RequestPayload& initialRequest,
                                         int ackCollectionTimeMs,
                                         int timeoutMs,
                                         CollaborationAckClosedHandler onAckClosed,
                                         ResponseHandler onFinalResponse,
                                         TimeoutHandler onTimeout,
                                         const RequestId& requestId = RequestId());
\end{Verbatim}
\textbf{源码：}\path{ndn-service-framework/ServiceUser.hpp}，第 858 行。\par

\Needspace{7\baselineskip}
\textbf{API-3ab0d3d87607}\quad public\par
\textbf{所属符号：}\code{ndn_service_framework::ServiceUser::CommitCollaborationPlan}\par
\begin{Verbatim}[fontsize=\footnotesize,breaklines,breakanywhere]
bool CommitCollaborationPlan(const RequestId& requestId,
                                         const std::string& ackClosedDigest,
                                         CollaborationPlan plan);
\end{Verbatim}
\textbf{源码：}\path{ndn-service-framework/ServiceUser.hpp}，第 887 行。\par

完整声明查询：\path{api/target-core-reference.md}。

\Needspace{12\baselineskip}
\section{Core API：调用流的事件和终态}
\lead{AC-06 · Core 接口契约}

\subsection{如何取得 writer}
User 通过 RequestServiceStreaming 建立单次调用并保存共享 handle；Provider 通过 addStreamingHandler 注册流处理器，由该处理器获得 StreamedResponseWriter。不要自行构造脱离 Core 请求绑定的 writer core。onEvent 接收 EventT，onComplete 接收最终 ResponseT，onError 接收 StreamedInvocationError；选定的模板类型需要可序列化。

\subsection{返回值与异常}

\begin{longtable}{@{}p{43mm}p{119mm}@{}}
\toprule
项目 & 行为与调用含义 \\
\midrule\endhead
publish(event) → bool & 模板先序列化再调用 core；任何异常被捕获并返回 false。false 也可能是 core 拒绝发布，不能仅解释为队列已满，更不能等同远端消费完成。 \\
finish(response, reason) & 返回 void；core 拒绝时抛 logic\_error（stream finish rejected）；序列化异常也可传播。不能在 finish 后继续发布普通事件。 \\
fail(code, message) & 错误终态已被占用时抛 logic\_error；与 publish 的吞异常返回 false 行为不同。 \\
取消/期限 & isCancelled 与 remainingDeadline 用于工作边界检查，仍需处理检查后发生取消的竞争；最终状态由 Core 仲裁。 \\
有界流 & options 管理期限、排队、保留与重试。事件留存可重传不等于业务执行可重复；应用不应通过再次执行生成来补一个遗失的网络事件。 \\
\bottomrule
\end{longtable}

\subsection{两事件示意}

\begin{Verbatim}[fontsize=\footnotesize,breaklines,breakanywhere]
// 在已认证的 Provider streaming handler 中
if (!writer.publish(event1)) { /* 停止产生新事件并检查关闭原因 */ }
else if (!writer.publish(event2)) { /* 同上 */ }
else { writer.finish(finalResponse, finishReason); }
// 必须处理 finish 的异常；不能假设 publish(false) 从不发生
\end{Verbatim}

\subsection{完整签名与源码定位}

\Needspace{7\baselineskip}
\textbf{API-9e9c476a768e}\quad public\par
\textbf{所属符号：}\code{ndn_service_framework::ServiceUser::RequestServiceStreaming}\par
\begin{Verbatim}[fontsize=\footnotesize,breaklines,breakanywhere]
template<typename RequestT, typename EventT, typename ResponseT>
std::shared_ptr<StreamedInvocationHandle<EventT, ResponseT>>
            RequestServiceStreaming(
                const ndn::Name& serviceName,
                const RequestT& request,
                StreamedInvocationOptions options,
                std::function<void(const EventT&)> onEvent,
                std::function<void(const ResponseT&)> onComplete,
                std::function<void(const StreamedInvocationError&)> onError,
                size_t strategy = ndn_service_framework::tlv::FirstResponding)
\end{Verbatim}
\textbf{源码：}\path{ndn-service-framework/ServiceUser.hpp}，第 779 行。\par

\Needspace{7\baselineskip}
\textbf{API-2eecc93af394}\quad public\par
\textbf{所属符号：}\code{ndn_service_framework::StreamedResponseWriter::publish}\par
\begin{Verbatim}[fontsize=\footnotesize,breaklines,breakanywhere]
bool publish(const EventT& event)
\end{Verbatim}
\textbf{源码：}\path{ndn-service-framework/InvocationStream.hpp}，第 348 行。\par

\Needspace{7\baselineskip}
\textbf{API-2bdd4eb83edb}\quad public\par
\textbf{所属符号：}\code{ndn_service_framework::StreamedResponseWriter::finish}\par
\begin{Verbatim}[fontsize=\footnotesize,breaklines,breakanywhere]
void finish(const ResponseT& response, StreamFinishReason reason)
\end{Verbatim}
\textbf{源码：}\path{ndn-service-framework/InvocationStream.hpp}，第 359 行。\par

完整声明查询：\path{api/target-core-reference.md}。

\Needspace{12\baselineskip}
\section{Core API：命名对象与连续流}
\lead{AC-07 · Core 接口契约}

\subsection{选择哪一种 API}

\begin{longtable}{@{}p{43mm}p{119mm}@{}}
\toprule
项目 & 行为与调用含义 \\
\midrule\endhead
签名对象 & publishSignedAppData 发布 dataName 对应的字节，freshness 是毫秒；fetchSignedAppData 传精确名字、expectedSigner、timeoutMs 和成功/失败回调。验证签名者后才能消费。 \\
完整大对象 & \code{publishEncryptedLargeData/fetchEncryptedLargeData} 或协作中的 \code{publishLargeNamed/fetchLargeExact} 处理分段与完整对象。manifest 引用和实际对象摘要由应用解释。 \\
连续流 & StreamPublisher::push 接收 shared\_ptr<ndn::Data> signedData，调用前已构造并签名；它不是接收未签名业务块并自动解释视频 metadata 的接口。 \\
调用事件流 & AC-06 的 writer 接收 EventT，由 Core 绑定单次请求、序号和终态。不能把它当成无结束时间的摄像头会话。 \\
\bottomrule
\end{longtable}

\subsection{连续流例子}

\begin{Verbatim}[fontsize=\footnotesize,breaklines,breakanywhere]
// signedData 已按流的精确名称/会话规则构造并签名
publisher.push(signedData); // void，不返回逐包送达确认
// subscriber 按会话与 cursor 获取并验证
// 应用决定过期视频帧可丢弃，完整录像对象必须完整验证
\end{Verbatim}

\subsection{资源与恢复}
发布窗口、订阅缓存和恢复范围有界，FEC 修复的是不透明字节，不能代替媒体解码或完整对象校验。发布、签名验证、重排、业务消费是不同阶段。调用者保存异步回调依赖的资源；流 stop 后历史 Data 仍可能在缓存中，应用必须用会话身份拒绝迟到旧帧。

\subsection{完整签名与源码定位}

\Needspace{7\baselineskip}
\textbf{API-160499aab055}\quad public\par
\textbf{所属符号：}\code{ndn_service_framework::ServiceUser::publishSignedAppData}\par
\begin{Verbatim}[fontsize=\footnotesize,breaklines,breakanywhere]
ndn::Name publishSignedAppData(
                const ndn::Name& dataName,
                const ndn::Buffer& payload,
                ndn::time::milliseconds freshness = ndn::DEFAULT_FRESHNESS_PERIOD);
\end{Verbatim}
\textbf{源码：}\path{ndn-service-framework/ServiceUser.hpp}，第 637 行。\par

\Needspace{7\baselineskip}
\textbf{API-ef2281d2fe2d}\quad public\par
\textbf{所属符号：}\code{ndn_service_framework::ServiceUser::fetchSignedAppData}\par
\begin{Verbatim}[fontsize=\footnotesize,breaklines,breakanywhere]
void fetchSignedAppData(
                const ndn::Name& dataName,
                const ndn::Name& expectedSigner,
                int timeoutMs,
                SignedAppDataHandler onData,
                SignedAppDataFailureHandler onFailure);
\end{Verbatim}
\textbf{源码：}\path{ndn-service-framework/ServiceUser.hpp}，第 660 行。\par

\Needspace{7\baselineskip}
\textbf{API-68706878d1a8}\quad public\par
\textbf{所属符号：}\code{ndn_service_framework::StreamPublisher::push}\par
\begin{Verbatim}[fontsize=\footnotesize,breaklines,breakanywhere]
void push(std::shared_ptr<ndn::Data> signedData);
\end{Verbatim}
\textbf{源码：}\path{ndn-service-framework/StreamFacade.hpp}，第 46 行。\par

完整声明查询：\path{api/target-core-reference.md}。

\Needspace{12\baselineskip}
\section{Core API：在线授权与版本状态}
\lead{AC-08 · Core 接口契约}

\subsection{输入是单项授权}
grant(identity, serviceName, authorizationAttribute) 接受一次身份/服务/属性操作。控制器把服务加入已有 allowedServices，把属性加入集合，再从有效属性重建完整 ABE 策略；调用者不提交该身份的全量策略。grant-only 推进 ControllerVersion，通常不旋转 ABE generation；如已有 pending rotation，先协调恢复，因此不能保证每次 grant 调用都绝无全局轮换。

\subsection{返回分支}

\begin{longtable}{@{}p{43mm}p{119mm}@{}}
\toprule
项目 & 行为与调用含义 \\
\midrule\endhead
grant 返回 false & 可能是非法 target、generation 未就绪、待轮换恢复失败、证书已撤销，或相同授权已经存在且无需重授的 no-op。bool 不能唯一编码失败原因。 \\
grant 产生变更 & 更新目标身份的权限与属性，移除相应服务撤销/可恢复的身份撤销；证书撤销不能靠 grant 覆盖。 \\
revoke(target) & RevocationTarget 显式区分身份、服务授权和证书。撤销事实/版本持久化后再做 ABE 轮换；后续失败不意味着撤销自动回滚。 \\
\code{getPolicyStatus(service)} & 得到服务相关撤销信息与全局版本；发布的外层 Data 要签名。policyDigest 不能当全部动态权限集合的规范摘要。 \\
\bottomrule
\end{longtable}

\subsection{接收者何时更新}
User/Provider 验证并安装 Controller 状态后，版本推进会使服务相关缓存失效；权限属性不变可以避免 DKEY 重取，但不保证业务状态毫无刷新。\code{NDNSF_POLICY_REVALIDATION_PERIOD_MS} 默认 0，未开启周期重校验时不能靠该机制主动发现。当前没有通用 SVS epoch 广播保证；认证 hint 仅触发权威拉取。目标 TG-01 将细化影响范围，不能用目标推定当前行为。

\subsection{完整签名与源码定位}

\Needspace{7\baselineskip}
\textbf{API-d5f6b3c470a6}\quad public\par
\textbf{所属符号：}\code{ndn_service_framework::ServiceController::grant}\par
\begin{Verbatim}[fontsize=\footnotesize,breaklines,breakanywhere]
bool grant(const ndn::Name& identity,
             const ndn::Name& serviceName,
             const ndn::Name& authorizationAttribute);
\end{Verbatim}
\textbf{源码：}\path{ndn-service-framework/ServiceController.hpp}，第 103 行。\par

\Needspace{7\baselineskip}
\textbf{API-3ad7bdf59cc1}\quad public\par
\textbf{所属符号：}\code{ndn_service_framework::ServiceController::revoke}\par
\begin{Verbatim}[fontsize=\footnotesize,breaklines,breakanywhere]
bool revoke(const RevocationTarget& target);
\end{Verbatim}
\textbf{源码：}\path{ndn-service-framework/ServiceController.hpp}，第 95 行。\par

\Needspace{7\baselineskip}
\textbf{API-a6e1ee327414}\quad public\par
\textbf{所属符号：}\code{ndn_service_framework::ServiceController::getPolicyStatus}\par
\begin{Verbatim}[fontsize=\footnotesize,breaklines,breakanywhere]
PolicyStatusData getPolicyStatus(const ndn::Name& serviceName) const;
\end{Verbatim}
\textbf{源码：}\path{ndn-service-framework/ServiceController.hpp}，第 91 行。\par

完整声明查询：\path{api/target-core-reference.md}。

\Needspace{12\baselineskip}
\section{Core API：执行租约、观测与并发}
\lead{AC-09 · Core 接口契约}

\subsection{租约携带什么}
GenericExecutionLease 包含 leaseId、\code{providerName/providerEpoch}、\code{requesterName/requestId/serviceName/planDigest}、\code{resourceBindingSchema/proof}、conflictKeys、expiresAtMs、executionDeadlineMs 与 idempotencyKey。providerEpoch 区分 Provider 进程代次；不能用 ControllerVersion 代替。resource proof 绑定资源，服务权限和 DI 解密 grant 另行检查。

\subsection{正常状态链}
prepare(lease, nowMs) 建立 Prepared；commit 用相同租约、Provider 代次、Requester 和幂等键进入 Committed；validateAndActivate 额外验证完整 ExecutionLeaseBinding 与执行期限，进入 Executing；结束后 release。未采用的预留可 abort，到期由 cleanupExpired 回收。find 返回 optional：查不到与找到了 Expired/Released 记录需要分别处理。

\subsection{结果与诊断}
ExecutionLeaseResult 包含 status、operation、reasonCode、lease、retryAfterMs、idempotentReplay。调用者先看 status，再按 reasonCode 处理 LEASE\_EXPIRED、LEASE\_STALE\_EPOCH、\code{LEASE_IDEMPOTENCY_CONFLICT}、\code{LEASE_REQUESTER_MISMATCH}、LEASE\_PLAN\_MISMATCH 或 LEASE\_CAPACITY\_REJECTED 等；重复成功回放由 idempotentReplay 标识。返回旧租约不等于可继续执行。

\subsection{并发与观察}
表内部以 mutex 保护状态，不使调用它的所有 Face API 任意线程安全。ACK/handler 工作池各有配置，增加线程不能扩大物理容量。counters(nowMs)、生命周期事件和 diagnostics 用于解释 \code{Prepared/Committed/Executing} 等资源占用；最终完成仍需实际结果证据。

\subsection{完整签名与源码定位}

\Needspace{7\baselineskip}
\textbf{API-d154a2bbea78}\quad public\par
\textbf{所属符号：}\code{ndn_service_framework::ProviderExecutionLeaseTable::prepare}\par
\begin{Verbatim}[fontsize=\footnotesize,breaklines,breakanywhere]
ExecutionLeaseResult
  prepare(GenericExecutionLease lease, uint64_t nowMs);
\end{Verbatim}
\textbf{源码：}\path{ndn-service-framework/ExecutionLease.hpp}，第 112 行。\par

\Needspace{7\baselineskip}
\textbf{API-13c34e09371a}\quad public\par
\textbf{所属符号：}\code{ndn_service_framework::ProviderExecutionLeaseTable::commit}\par
\begin{Verbatim}[fontsize=\footnotesize,breaklines,breakanywhere]
ExecutionLeaseResult
  commit(const std::string& leaseId, const std::string& providerEpoch,
         const std::string& requesterName, const std::string& idempotencyKey,
         uint64_t nowMs);
\end{Verbatim}
\textbf{源码：}\path{ndn-service-framework/ExecutionLease.hpp}，第 115 行。\par

\Needspace{7\baselineskip}
\textbf{API-e156fa6d6cf2}\quad public\par
\textbf{所属符号：}\code{ndn_service_framework::ProviderExecutionLeaseTable::validateAndActivate}\par
\begin{Verbatim}[fontsize=\footnotesize,breaklines,breakanywhere]
ExecutionLeaseResult
  validateAndActivate(const std::string& leaseId,
                      const std::string& providerEpoch,
                      const ExecutionLeaseBinding& binding,
                      const std::string& idempotencyKey,
                      uint64_t nowMs,
                      uint64_t executionDeadlineMs);
\end{Verbatim}
\textbf{源码：}\path{ndn-service-framework/ExecutionLease.hpp}，第 120 行。\par

\Needspace{7\baselineskip}
\textbf{API-f52dcb87cd76}\quad public\par
\textbf{所属符号：}\code{ndn_service_framework::ProviderExecutionLeaseTable::release}\par
\begin{Verbatim}[fontsize=\footnotesize,breaklines,breakanywhere]
ExecutionLeaseResult
  release(const std::string& leaseId, const std::string& providerEpoch,
          const std::string& requesterName, const std::string& idempotencyKey,
          uint64_t nowMs);
\end{Verbatim}
\textbf{源码：}\path{ndn-service-framework/ExecutionLease.hpp}，第 145 行。\par

\Needspace{7\baselineskip}
\textbf{API-e43c54567dc6}\quad public\par
\textbf{所属符号：}\code{ndn_service_framework::ProviderExecutionLeaseTable::find}\par
\begin{Verbatim}[fontsize=\footnotesize,breaklines,breakanywhere]
std::optional<GenericExecutionLease>
  find(const std::string& leaseId) const;
\end{Verbatim}
\textbf{源码：}\path{ndn-service-framework/ExecutionLease.hpp}，第 153 行。\par

完整声明查询：\path{api/target-core-reference.md}。

\Needspace{12\baselineskip}
\section{Repo API：对象与精确 Data packet}
\lead{AC-10 · Repo 接口契约}

\subsection{对象与 manifest}
put(node, objectName, payload, options) 的本地重载接收 RepoNode 和字节，返回 RepoObjectManifest。StoreOptions 提供 objectType、replicationFactor、replicaNodes、policyEpoch。名字识别对象，摘要识别内容，segmentCount 描述分段数量，副本字段只是位置/策略信息，不能用其中一个替代全部验证。网络 helper 与本地静态调用不是同一传输或超时模型。

\subsection{原始 packet 例子}
录像含两份已经由摄像头签名的 Data，名称分别指向两个精确 packet。putDataPacket 保存原始 wire；manifest 保存有序完整名称；getDataPackets 按该列表读回。必须核对完整组的名称和索引，不能重新签名 payload 来冒充原摄像头。普通字节对象 put 则不承诺这些字节本身已经加密。

\subsection{存储与可见性}
普通对象由 SQLite 权威存储加有界缓存支持，权威写入后更新缓存；RepoCore 负责网络中立操作，RepoNode 持有网络服务注册。remove 的本地删除结果不能表示远端副本、目录或 NDN 缓存同时消失。取回 manifest 只是下一步取数依据，完整对象成功须验证全部字节及声明摘要。

\subsection{完整签名与源码定位}

\Needspace{7\baselineskip}
\textbf{API-3690bad70c85}\quad public\par
\textbf{所属符号：}\code{ndnsf_distributed_repo::RepoClient::put}\par
\begin{Verbatim}[fontsize=\footnotesize,breaklines,breakanywhere]
static RepoObjectManifest put(RepoNode& node,
                                const std::string& objectName,
                                const std::vector<uint8_t>& payload,
                                StoreOptions options = {});
\end{Verbatim}
\textbf{源码：}\path{NDNSF-DistributedRepo/include/ndnsf-distributed-repo/RepoClient.hpp}，第 33 行。\par

\Needspace{7\baselineskip}
\textbf{API-208878700957}\quad public\par
\textbf{所属符号：}\code{ndnsf_distributed_repo::RepoClient::getDataPackets}\par
\begin{Verbatim}[fontsize=\footnotesize,breaklines,breakanywhere]
static std::vector<std::vector<uint8_t>> getDataPackets(
    const RepoNode& node,
    const RepoObjectManifest& manifest);
\end{Verbatim}
\textbf{源码：}\path{NDNSF-DistributedRepo/include/ndnsf-distributed-repo/RepoClient.hpp}，第 62 行。\par

完整声明查询：\path{api/target-repo-reference.md}。

\Needspace{12\baselineskip}
\section{Repo API：制品上传与提交}
\lead{AC-11 · Repo 接口契约}

\subsection{公开上传调用}
publish\_file 要求已有文件 path、逻辑 name、expected\_sha256；replicas 默认 1、verification 默认 signed-manifest、resume 默认 True。应用必须先独立得到预期内容摘要，不能从未验证的上传结果倒填。timeout\_ms、取消 token 与 progress observer 控制不同方面。

\subsection{显式阶段}

\begin{Verbatim}[fontsize=\footnotesize,breaklines,breakanywhere]
// Python 调用顺序示意
session = api.begin_upload(descriptor)
session.upload_file(path)
result = session.commit()
// transfer/upload 成功仍不能替代 commit 成功
// 失败时按恢复策略 abort(preserve_progress=...)
\end{Verbatim}

\subsection{结果字段}

\begin{longtable}{@{}p{43mm}p{119mm}@{}}
\toprule
项目 & 行为与调用含义 \\
\midrule\endhead
reference / operation\_id & reference 是后续下载依据，operation\_id 识别本次操作；不能仅保存本地文件路径。 \\
requested / achieved replicas & 请求数量与成功数量分别报告；achieved\_replicas 必须等于不同 COMMITTED receipt\_id 的数量，不能只数候选 ACK。 \\
replicas & 每个 ArtifactReplicaResult 含 repo\_node、state、receipt\_id、error\_code；用于识别哪个副本完成，哪个失败。 \\
deduplicated / resumed & 分别说明复用与恢复，不代表本轮传输了全部字节；duration 和 phase\_durations\_ms 是本次观测。 \\
\bottomrule
\end{longtable}

\subsection{提交与重试}
commit 校验 operation、摘要、选定副本及 receipt 后产生可交付结果。幂等键用于识别同一业务上传，但不能允许同一操作携带冲突内容。网络后端的 TARGETED 全制品发布支持有边界，不能由控制选项存在推断所有模式可用。失败保留 operation 和阶段原因；是否保留进度按 resume/abort 策略决定。

\subsection{完整签名与源码定位}

\Needspace{7\baselineskip}
\textbf{API-c7c484fc1ed7}\quad public-by-name\par
\textbf{所属符号：}\code{ArtifactRepositoryApi.publish_file}\par
\begin{Verbatim}[fontsize=\footnotesize,breaklines,breakanywhere]
def publish_file(
        self,
        path: Union[str, Path],
        *,
        name: str,
        expected_sha256: str,
        replicas: int = 1,
        verification: str = "signed-manifest",
        resume: bool = True,
        on_progress: Optional[ProgressObserver] = None,
        idempotency_key: str = "",
        policy_epoch: str = "default",
        timeout_ms: Optional[int] = None,
        control: ArtifactControlOptions = ArtifactControlOptions(),
        cancellation: Optional[ArtifactCancellationToken] = None,
    ) -> ArtifactPublishResult:
\end{Verbatim}
\textbf{源码：}\path{NDNSF-DistributedRepo/pythonWrapper/py_repoclient/artifact_api.py}，第 844 行。\par

\Needspace{7\baselineskip}
\textbf{API-91dcd6af85dd}\quad public-by-name\par
\textbf{所属符号：}\code{ArtifactRepositoryApi.begin_upload}\par
\begin{Verbatim}[fontsize=\footnotesize,breaklines,breakanywhere]
def begin_upload(
        self,
        descriptor: ArtifactDescriptor,
        *,
        on_progress: Optional[ProgressObserver] = None,
        cancellation: Optional[ArtifactCancellationToken] = None,
    ) -> ArtifactUploadSession:
\end{Verbatim}
\textbf{源码：}\path{NDNSF-DistributedRepo/pythonWrapper/py_repoclient/artifact_api.py}，第 819 行。\par

完整声明查询：\path{api/target-repo-reference.md}。

\Needspace{12\baselineskip}
\section{Repo API：下载、验证与传输状态机}
\lead{AC-12 · Repo 接口契约}

\subsection{kwargs 的实际含义}
fetch\_file(reference, destination, **kwargs) 把参数传给 begin\_fetch：resume=True、verify=True、replace=False、on\_progress=None、idempotency\_key 空、timeout\_ms=None、control=\code{ArtifactControlOptions()}、cancellation=None。默认幂等键由 \code{reference.content_digest} 与目标绝对路径构造，operation identity 还绑定 reference；换路径并不天然延续同一操作。timeout\_ms 通过 int(timeout\_ms or default\_timeout\_ms) 传给后端，所以传 0 会采用默认预算。

\subsection{调用次序与结果}

\begin{Verbatim}[fontsize=\footnotesize,breaklines,breakanywhere]
result = api.fetch_file(reference, destination,
                        resume=True, verify=True, replace=False)
// 等价主线：begin_fetch → transfer → commit
// 异常且尚未终态：尝试 abort(preserve_progress=resume)，再抛 API error
\end{Verbatim}

\subsection{如何读统计}
ArtifactFetchResult 包含 reference、operation\_id、destination、reused\_bytes、transferred\_bytes、source\_replicas、last\_segment、delivered\_segments、total\_segments、retransmitted\_bytes 和阶段耗时。复用字节与本轮网络字节分开；last\_segment 默认 -1，不能当作完整对象成功。成功交付来自 commit 和校验，不由进度达到某个百分比推断。

\subsection{底层分段是另一层状态}
\code{AdaptiveArtifactTransfer::receive} 记录 segmentNo、logicalBytes、wireBytes 和 nowMs，返回 \code{ArtifactSegmentDisposition}；markVerified 显式标记验证完成。接收、验证、整体提交是不同动作。高层 fetch session 管理文件、操作身份和异常恢复，底层窗口/重传状态不能直接当成最终文件完成。resume 必须保留同一制品和可信分段状态，不是无条件追加目标文件。

\subsection{完整签名与源码定位}

\Needspace{7\baselineskip}
\textbf{API-39b8d985de2c}\quad public-by-name\par
\textbf{所属符号：}\code{ArtifactRepositoryApi.begin_fetch}\par
\begin{Verbatim}[fontsize=\footnotesize,breaklines,breakanywhere]
def begin_fetch(
        self,
        reference: ArtifactReference,
        destination: Union[str, Path],
        *,
        resume: bool = True,
        verify: bool = True,
        replace: bool = False,
        on_progress: Optional[ProgressObserver] = None,
        idempotency_key: str = "",
        timeout_ms: Optional[int] = None,
        control: ArtifactControlOptions = ArtifactControlOptions(),
        cancellation: Optional[ArtifactCancellationToken] = None,
    ) -> ArtifactFetchSession:
\end{Verbatim}
\textbf{源码：}\path{NDNSF-DistributedRepo/pythonWrapper/py_repoclient/artifact_api.py}，第 914 行。\par

\Needspace{7\baselineskip}
\textbf{API-1707faa2f3ea}\quad public-by-name\par
\textbf{所属符号：}\code{ArtifactRepositoryApi.fetch_file}\par
\begin{Verbatim}[fontsize=\footnotesize,breaklines,breakanywhere]
def fetch_file(
        self,
        reference: ArtifactReference,
        destination: Union[str, Path],
        **kwargs,
    ) -> ArtifactFetchResult:
\end{Verbatim}
\textbf{源码：}\path{NDNSF-DistributedRepo/pythonWrapper/py_repoclient/artifact_api.py}，第 960 行。\par

\Needspace{7\baselineskip}
\textbf{API-6f363704fd88}\quad public\par
\textbf{所属符号：}\code{ndnsf_distributed_repo::AdaptiveArtifactTransfer::receive}\par
\begin{Verbatim}[fontsize=\footnotesize,breaklines,breakanywhere]
ArtifactSegmentDisposition
  receive(uint64_t segmentNo, uint64_t logicalBytes, uint64_t wireBytes,
          uint64_t nowMs);
\end{Verbatim}
\textbf{源码：}\path{NDNSF-DistributedRepo/include/ndnsf-distributed-repo/ArtifactTransfer.hpp}，第 78 行。\par

\Needspace{7\baselineskip}
\textbf{API-b8a723742c4e}\quad public\par
\textbf{所属符号：}\code{ndnsf_distributed_repo::AdaptiveArtifactTransfer::markVerified}\par
\begin{Verbatim}[fontsize=\footnotesize,breaklines,breakanywhere]
void markVerified(uint64_t segmentNo);
\end{Verbatim}
\textbf{源码：}\path{NDNSF-DistributedRepo/include/ndnsf-distributed-repo/ArtifactTransfer.hpp}，第 82 行。\par

完整声明查询：\path{api/target-repo-reference.md}。

\Needspace{12\baselineskip}
\section{Repo API：目录、副本与缓存控制}
\lead{AC-13 · Repo 接口契约}

\subsection{本地目录操作矩阵}

\begin{longtable}{@{}p{43mm}p{119mm}@{}}
\toprule
项目 & 行为与调用含义 \\
\midrule\endhead
catalogStatus & 本 Repo 的模式、catalogEpoch 和对象数；不是所有 Repo 的一致全局视图。 \\
catalogSnapshot & 返回 RepoCatalogDelta 形式的全量目录；调用方按快照恢复自己的视图。 \\
\code{catalogDelta(sinceEpoch)} & 按本 Repo 的目录 epoch 请求变化；空变化与查不到单个对象不是同一个结果，不能把其他 Repo 的 epoch 当作游标。 \\
\code{catalogLookup(objectName)} & 精确读取一个对象的 manifest，再生成 AVAILABLE 条目；无类型、类别、前缀或标签过滤参数。 \\
Python orchestration & 查询编排、目录交换和 repair 属于独立路径；C++ 类型存在不表示已在 C++ 启动这些服务。 \\
\bottomrule
\end{longtable}

\subsection{lookup 的实际失败与锁边界}
RepoCore 先在目录锁外调用 store->get(objectName)，再锁住目录状态构造条目和 catalogEpoch。对象缺失或存储异常沿路径传播，不承诺返回空集合。该顺序不保证存储与目录是一份线性一致快照；并发删除/更新时，调用方不可将条目视为永久可读凭据。handleCatalogLookup 将请求字节解释为对象名，返回条目 JSON。

\subsection{副本判断例子}
目录声称 P1、P2 保存 artifact X 时，下载还需验证 X 的 manifest 和实际字节；若 P2 不可读，不能继续统计为本轮已验证的第二副本。repair 应遵守 tombstone、retention 和 repairAllowed；旧目录 AVAILABLE 信息不能无条件复活已删除对象。TG-04 将进一步规定跨后端能力及提交/恢复一致性。

\subsection{完整签名与源码定位}

\Needspace{7\baselineskip}
\textbf{API-b28b415568be}\quad public\par
\textbf{所属符号：}\code{ndnsf_distributed_repo::RepoCore::catalogSnapshot}\par
\begin{Verbatim}[fontsize=\footnotesize,breaklines,breakanywhere]
RepoCatalogDelta catalogSnapshot() const;
\end{Verbatim}
\textbf{源码：}\path{NDNSF-DistributedRepo/include/ndnsf-distributed-repo/RepoCore.hpp}，第 60 行。\par

\Needspace{7\baselineskip}
\textbf{API-e296daab62b3}\quad public\par
\textbf{所属符号：}\code{ndnsf_distributed_repo::RepoCore::catalogDelta}\par
\begin{Verbatim}[fontsize=\footnotesize,breaklines,breakanywhere]
RepoCatalogDelta catalogDelta(uint64_t sinceEpoch) const;
\end{Verbatim}
\textbf{源码：}\path{NDNSF-DistributedRepo/include/ndnsf-distributed-repo/RepoCore.hpp}，第 62 行。\par

\Needspace{7\baselineskip}
\textbf{API-2c18852cdee0}\quad public\par
\textbf{所属符号：}\code{ndnsf_distributed_repo::RepoCore::catalogLookup}\par
\begin{Verbatim}[fontsize=\footnotesize,breaklines,breakanywhere]
RepoCatalogEntry catalogLookup(const std::string& objectName) const;
\end{Verbatim}
\textbf{源码：}\path{NDNSF-DistributedRepo/include/ndnsf-distributed-repo/RepoCore.hpp}，第 64 行。\par

完整声明查询：\path{api/target-repo-reference.md}。

\Needspace{12\baselineskip}
\section{DI API：Python 应用入口与原生句柄}
\lead{AC-14 · DI 接口契约}

\subsection{Python 与 C++ 是两个现存入口}
APPClient.request 输入 model、task、input、timeout\_ms，可带 \code{options/objective/constraints/strategy}。model 描述模型身份，task 指定业务任务，input 是应用值，策略决定候选选择；显式 deployment 的 request 入口属于另一种调用形态。Python api 包重新导出 app\_sdk，不意味着运行时已经委托给 C++。

\subsection{原生输入与选项}
NativeApplicationInput 区分 Inline 和 RepositoryReference，携带任务及 schema 绑定；内联字节与 Repo 引用不能混填。NativeRequestOptions 的 timeoutMs 默认 30000，ackTimeoutMs 默认 5000，后者是前者之内的报价预算。NativeInferenceClient 创建 handle 并调度请求；目前生产 dispatch 仍在期限检查后返回 \code{NATIVE_REQUEST_PIPELINE_NOT_READY}。

\subsection{句柄行为}

\begin{longtable}{@{}p{43mm}p{119mm}@{}}
\toprule
项目 & 行为与调用含义 \\
\midrule\endhead
result(waitTimeout) & 局部等待超时抛 LOCAL\_WAIT\_TIMEOUT，请求仍可继续；不修改请求 deadline。Core I/O 线程正数等待预算抛 CORE\_IO\_WAIT\_FORBIDDEN。 \\
cancel & 空句柄或已终态请求不重新提交终态；返回并不保证远端计算已经停止，也不等价于释放 Provider KV。 \\
observe & 拒绝空 observer，回放已有事件；回放与新事件在独立通知队列串行排序。回调异常被捕获，捕获资源须存活至通知结束。 \\
close / 生命周期 & client close 取消当前操作并关闭调度资源；等待者持有共享 operation。迟到 dispatch 在 phase/状态不符时退出，不复活请求。 \\
NativeDiError & 保留 \code{code/domain/boundary/requestId/attempt}；调用者按错误字段分流，不能仅比较一段异常文案。 \\
\bottomrule
\end{longtable}

\subsection{等待例子}
应用给请求 30 秒预算，先用 100 毫秒局部等待显示进度。出现 LOCAL\_WAIT\_TIMEOUT 后可继续观察或再次等待，不应创建重复推理请求；只有应用明确结束业务才 cancel。这个例子解释句柄语义，不代表当前原生 requester 已能成功完成模型推理；接通路线见 TG-02。

\subsection{完整签名与源码定位}

\Needspace{7\baselineskip}
\textbf{API-0ad867060005}\quad public-by-name\par
\textbf{所属符号：}APPClient.request\par
\begin{Verbatim}[fontsize=\footnotesize,breaklines,breakanywhere]
def request(
        self,
        *,
        model,
        task,
        input,
        timeout_ms: int,
        options=None,
        objective=None,
        constraints=None,
        request_id: str = "",
        strategy=None,
    ):
\end{Verbatim}
\textbf{源码：}\path{NDNSF-DistributedInference/ndnsf_distributed_inference/app_sdk/client.py}，第 307 行。\par

\Needspace{7\baselineskip}
\textbf{API-375b92c9a402}\quad public\par
\textbf{所属符号：}\code{ndnsf::di::NativeInferenceClient::request}\par
\begin{Verbatim}[fontsize=\footnotesize,breaklines,breakanywhere]
NativeInferenceHandle request(
    const NativeModelRef& model,
    const NativeApplicationInput& input,
    std::shared_ptr<const NativeModelSplitStrategy> splitStrategy,
    std::shared_ptr<const NativePlacementStrategy> placementStrategy,
    const NativeRequestOptions& options);
\end{Verbatim}
\textbf{源码：}\path{NDNSF-DistributedInference/cpp/ndnsf-di/NativeInferenceClient.hpp}，第 121 行。\par

\Needspace{7\baselineskip}
\textbf{API-c0c25988d534}\quad public\par
\textbf{所属符号：}\code{ndnsf::di::NativeInferenceHandle::result}\par
\begin{Verbatim}[fontsize=\footnotesize,breaklines,breakanywhere]
NativeInferenceResult result(std::chrono::milliseconds waitTimeout) const;
\end{Verbatim}
\textbf{源码：}\path{NDNSF-DistributedInference/cpp/ndnsf-di/NativeInferenceClient.hpp}，第 95 行。\par

完整声明查询：\path{api/target-di-reference.md}。

\Needspace{12\baselineskip}
\section{DI API：模型适配与角色放置扩展}
\lead{AC-15 · DI 接口契约}

\subsection{规划对象怎么连接}

\begin{longtable}{@{}p{43mm}p{119mm}@{}}
\toprule
项目 & 行为与调用含义 \\
\midrule\endhead
模型描述 / 图快照 & 描述模型及 adapter 身份；inspect 得到真实图，planning graph 与 canonical ONNX graph 可有不同摘要，不能只因都叫 graphDigest 就混用。 \\
NativeSplitCandidate & enumerate 返回的候选包含模型/图绑定和角色方案；候选数量受预算控制。候选是可验证的执行分解，不是任意 Provider 名单。 \\
ProviderOfferV3 & 来自认证 ACK 的报价，经 admission 检查请求、attempt、模型/图、Provider、设备和新鲜度；不是从 capability 文本自行补全的运行承诺。 \\
V3 proposal & proposeRoles 以绑定上下文、ackClosedDigest、角色、合格报价和 nowMs 提出放置；独立 validator 仍要校验覆盖、资源和身份。 \\
\bottomrule
\end{longtable}

\subsection{两段模型例子}
将线性模型的前半段作为 A、后半段作为 B，候选须说明 A 的输出与 B 的输入匹配。P1 报价能运行 A，P2 报价能运行 B，策略才可提出 A→P1、B→P2；如果 P2 只有模糊 has\_model 声明而缺该角色/设备证明，不能标为 exact reuse。若允许准备，则明确采用需装配/加载的路径，并把准备成本计入总期限。

\subsection{扩展最小职责}
\code{NativeModelSplitStrategy} 实现 identity 和 enumerate；NativePlacementStrategy 实现 identity 和 proposeRoles。adapter 拥有模型检查、输入编码、结果解码，由 registry 登记并在使用前固定。策略返回值不能给自己签发 grant。新增策略应以固定图、报价、时钟测试确定性和拒绝分支，再让实际 requester 选择它；仅创建子类不会改变默认路径。

\subsection{完整签名与源码定位}

\Needspace{7\baselineskip}
\textbf{API-da75fc521e8b}\quad public\par
\textbf{所属符号：}\code{ndnsf::di::NativeModelSplitStrategy::enumerate}\par
\begin{Verbatim}[fontsize=\footnotesize,breaklines,breakanywhere]
virtual std::vector<NativeSplitCandidate> enumerate(
    const NativeModelDescriptor& model,
    const NativeGraphSnapshot& graph,
    const NativeCandidateBudget& budget) const = 0;
\end{Verbatim}
\textbf{源码：}\path{NDNSF-DistributedInference/cpp/ndnsf-di/NativePlanning.hpp}，第 172 行。\par

\Needspace{7\baselineskip}
\textbf{API-0f4ac0e590fd}\quad public\par
\textbf{所属符号：}\code{ndnsf::di::NativePlacementStrategy::proposeRoles}\par
\begin{Verbatim}[fontsize=\footnotesize,breaklines,breakanywhere]
virtual NativeRolePlacementProposalV3 proposeRoles(
    const NativeOfferBindingContext& context, const std::string& ackClosedDigest,
    const std::vector<NativeSelectionRoleV3>& roles,
    const std::vector<NativeAdmittedOfferV3>& offers, std::uint64_t nowMs) const = 0;
\end{Verbatim}
\textbf{源码：}\path{NDNSF-DistributedInference/cpp/ndnsf-di/NativePlanning.hpp}，第 183 行。\par

完整声明查询：\path{api/target-di-reference.md}。

\Needspace{12\baselineskip}
\section{DI API：准备、发布后再绑定与封存}
\lead{AC-16 · DI 接口契约}

\subsection{阶段输入输出}

\begin{longtable}{@{}p{43mm}p{119mm}@{}}
\toprule
项目 & 行为与调用含义 \\
\midrule\endhead
prepareInput / inspectModel & 绑定 task/schema、输入字节或引用和 NativeRequestControl；检查真实模型与图。control 保存 request/attempt、steady-clock deadline 和取消检查。 \\
prepareRoles / validateRoles & 把候选转成可执行角色并检查一致性，不能先封存再猜测模型边界。 \\
ensureArtifacts & 消费完整 candidate、proposal、control，通过发布端口确保制品；检查业务 root 原始字节及模型/profile/source/package 等身份。 \\
bindPublishedRoles & 将发布后的 manifest/recipe 绑定到角色；保持角色、候选、Provider、设备选择不被偷偷更换，旧 exact-reuse 证明必须仍覆盖实际 recipe。 \\
sealCore / grantView & 前者固定计划核心，后者形成授权所需视图；摘要对应具体规范对象，不是任意 JSON 文本。 \\
finalizeSecurity / project / encode & 完成安全材料后产生各 Provider 的选择投影并编码；不能在安全材料未绑定时发送一个看似完整的旧投影。 \\
\bottomrule
\end{longtable}

\subsection{发布改变身份的例子}
候选最初引用 canonical 模型，发布后得到切片 artifact root R 和装配 recipe Q。不能只把 URL 改为 R 而继续保留旧 planDigest 或驻留证明；bindPublishedRoles 校验选择不变并绑定 R/Q，再封存对应计划与授权。业务 root 摘要标识模型制品，Core 传输 manifest 描述取数封装，两者不能互代。

\subsection{失败边界}
取消/过期先由 control 拒绝继续准备；内容或角色不匹配在准备/再绑定阶段拒绝，grant 验证失败在安全边界拒绝。端口定向测试不证明生产 Repo/catalog 与 requester 已接通；实现状态集中见第 31 章与 Spec182，而非给每一步追加成功承诺。

\subsection{完整签名与源码定位}

\Needspace{7\baselineskip}
\textbf{API-f8cbacddff79}\quad public\par
\textbf{所属符号：}\code{ndnsf::di::NativeRequestPreparation::prepareInput}\par
\begin{Verbatim}[fontsize=\footnotesize,breaklines,breakanywhere]
NativePreparedInput prepareInput(const NativeModelDescriptor& model,
                                   std::string taskName,
                                   std::string inputSchemaDigest,
                                   std::string optionsSchemaDigest,
                                   std::vector<std::uint8_t> payload,
                                   std::string repositoryReference,
                                   std::chrono::steady_clock::time_point deadline) const;
\end{Verbatim}
\textbf{源码：}\path{NDNSF-DistributedInference/cpp/ndnsf-di/NativeRequestPreparation.hpp}，第 105 行。\par

\Needspace{7\baselineskip}
\textbf{API-dcca376d355a}\quad public\par
\textbf{所属符号：}\code{ndnsf::di::NativeRequestPreparation::inspectModel}\par
\begin{Verbatim}[fontsize=\footnotesize,breaklines,breakanywhere]
NativeInspectedModel inspectModel(const NativePreparedInput& input) const;
\end{Verbatim}
\textbf{源码：}\path{NDNSF-DistributedInference/cpp/ndnsf-di/NativeRequestPreparation.hpp}，第 113 行。\par

\Needspace{7\baselineskip}
\textbf{API-2d8e39c7fc67}\quad public\par
\textbf{所属符号：}\code{ndnsf::di::NativeRequestPreparation::ensureArtifacts}\par
\begin{Verbatim}[fontsize=\footnotesize,breaklines,breakanywhere]
NativeArtifactBinding ensureArtifacts(const NativeInspectedModel& model,
                                        const NativeSplitCandidate& candidate,
                                        const NativeRolePlacementProposalV3& proposal,
                                        const NativeRequestControl& control) const;
\end{Verbatim}
\textbf{源码：}\path{NDNSF-DistributedInference/cpp/ndnsf-di/NativeRequestPreparation.hpp}，第 121 行。\par

\Needspace{7\baselineskip}
\textbf{API-d192a8a670da}\quad public\par
\textbf{所属符号：}\code{ndnsf::di::NativeRequestPreparation::bindPublishedRoles}\par
\begin{Verbatim}[fontsize=\footnotesize,breaklines,breakanywhere]
static std::vector<NativeSelectionRoleV3> bindPublishedRoles(
    const NativeInspectedModel& model, const NativeSplitCandidate& candidate,
    const std::vector<NativeSelectionRoleV3>& roles, const NativeArtifactBinding& artifacts);
\end{Verbatim}
\textbf{源码：}\path{NDNSF-DistributedInference/cpp/ndnsf-di/NativeRequestPreparation.hpp}，第 128 行。\par

\Needspace{7\baselineskip}
\textbf{API-1df2153e06a6}\quad public\par
\textbf{所属符号：}\code{ndnsf::di::NativePlanSealer::sealCore}\par
\begin{Verbatim}[fontsize=\footnotesize,breaklines,breakanywhere]
static NativePlacementPlanCore sealCore(
    const NativePlanningSnapshot& snapshot,
    const NativePlacementProposal& proposal,
    const NativePlanSealingInputs& inputs);
\end{Verbatim}
\textbf{源码：}\path{NDNSF-DistributedInference/cpp/ndnsf-di/NativePlanSealer.hpp}，第 114 行。\par

\Needspace{7\baselineskip}
\textbf{API-550c3b9b50b5}\quad public\par
\textbf{所属符号：}\code{ndnsf::di::NativePlanSealer::grantView}\par
\begin{Verbatim}[fontsize=\footnotesize,breaklines,breakanywhere]
static NativeProviderGrantView grantView(
    const NativePlacementPlanCore& core,
    const NativeProviderPlanningView& provider,
    const NativeSecurityPolicySnapshot& security);
\end{Verbatim}
\textbf{源码：}\path{NDNSF-DistributedInference/cpp/ndnsf-di/NativePlanSealer.hpp}，第 125 行。\par

\Needspace{7\baselineskip}
\textbf{API-e3e459f19b03}\quad public\par
\textbf{所属符号：}\code{ndnsf::di::NativePlanSealer::project}\par
\begin{Verbatim}[fontsize=\footnotesize,breaklines,breakanywhere]
static NativeSelectionProjectionV3 project(
    const NativeSealedPlan& sealed, const std::string& provider,
    const NativeRoleProjectionInputs& inputs);
\end{Verbatim}
\textbf{源码：}\path{NDNSF-DistributedInference/cpp/ndnsf-di/NativePlanSealer.hpp}，第 138 行。\par

\Needspace{7\baselineskip}
\textbf{API-07a74268daf4}\quad public\par
\textbf{所属符号：}\code{ndnsf::di::NativePlanSealer::encode}\par
\begin{Verbatim}[fontsize=\footnotesize,breaklines,breakanywhere]
static std::vector<std::uint8_t> encode(
    const NativeSelectionProjectionV3& projection);
\end{Verbatim}
\textbf{源码：}\path{NDNSF-DistributedInference/cpp/ndnsf-di/NativePlanSealer.hpp}，第 142 行。\par

完整声明查询：\path{api/target-di-reference.md}。

\Needspace{12\baselineskip}
\section{DI API：Provider 注册、授权与执行}
\lead{AC-17 · DI 接口契约}

\subsection{组成一个 Provider}
NativeServiceDefinition 声明 serviceName 和模型支持，\code{NativeProviderHandlerConfig} 提供处理配置；serve 返回必须保存的 \code{NativeServiceRegistration}。执行层 NativeProviderSession 持有 plan、assignment、DependencyIo、runnerFactory 和 worker 数，按角色 registerRunner。plan 描述做什么，assignment 指定谁做，DependencyIo 指定怎样获取/发布边数据，runnerFactory 提供实际后端。

\subsection{执行次序}
Core scoped handler 接受并验证 Selection 后，NativeProviderHandler 取得协作认证上下文，核对计划、角色、lease 和 grant；再构造/调用 session。executeRoleAsync 返回 future<ProviderRoleResult>，调用方须处理 future 的异常并保留 session/I/O 生命周期。source role 可用初始输入，有依赖的角色等直接前驱输出，不需要凭空增加全局所有节点 Ready 屏障。

\subsection{grant 从哪里来}
\code{verifyAndUnwrapNativeGrant} 的输入来自认证选择上下文和 Provider 投影，不能使用应用随意填写的 requester/plan 字符串代替已认证来源。服务权限允许调用，execution lease 约束资源，request-scoped grant 约束受保护模型与计划；三个条件不能互相代替。默认 plaintext-v1 路径与 ProtectedRuntime 明文生命周期管理要分别说明。

\subsection{关闭与迟到工作}
保存服务注册句柄直到停服，close 阻断该 generation 的新工作和晚到发布；应用还需排空已启动 worker。关闭服务、释放运行 lease、释放缓存状态和删除 artifact 是不同操作。真实后端/设备以及迁移接线状态由第 31 章与当前 Spec 记录。

\subsection{完整签名与源码定位}

\Needspace{7\baselineskip}
\textbf{API-b5abfc9b9043}\quad public\par
\textbf{所属符号：}\code{ndnsf::di::NativeInferenceProvider::serve}\par
\begin{Verbatim}[fontsize=\footnotesize,breaklines,breakanywhere]
NativeServiceRegistration serve(const NativeServiceDefinition& service,
                                  const NativeProviderHandlerConfig& config);
\end{Verbatim}
\textbf{源码：}\path{NDNSF-DistributedInference/cpp/ndnsf-di/NativeInferenceProvider.hpp}，第 98 行。\par

\Needspace{7\baselineskip}
\textbf{API-1ec58fcbe7cc}\quad public\par
\textbf{所属符号：}\code{ndnsf::di::NativeProviderSession::NativeProviderSession}\par
\begin{Verbatim}[fontsize=\footnotesize,breaklines,breakanywhere]
NativeProviderSession(NativeExecutionPlan plan,
                        NativeProviderAssignment assignment,
                        std::shared_ptr<DependencyIo> dependencyIo,
                        std::shared_ptr<NativeModelRunnerFactory> runnerFactory,
                        std::size_t workerCount = std::thread::hardware_concurrency());
\end{Verbatim}
\textbf{源码：}\path{NDNSF-DistributedInference/cpp/ndnsf-di/NativeProviderSession.hpp}，第 19 行。\par

\Needspace{7\baselineskip}
\textbf{API-79b6adfd9bf2}\quad public\par
\textbf{所属符号：}\code{ndnsf::di::NativeProviderSession::registerRunner}\par
\begin{Verbatim}[fontsize=\footnotesize,breaklines,breakanywhere]
void
  registerRunner(const NativeModelRunnerSpec& spec);
\end{Verbatim}
\textbf{源码：}\path{NDNSF-DistributedInference/cpp/ndnsf-di/NativeProviderSession.hpp}，第 25 行。\par

\Needspace{7\baselineskip}
\textbf{API-5ea3dea813e9}\quad public\par
\textbf{所属符号：}\code{ndnsf::di::NativeProviderSession::executeRoleAsync}\par
\begin{Verbatim}[fontsize=\footnotesize,breaklines,breakanywhere]
std::future<ProviderRoleResult>
  executeRoleAsync(const std::string& sessionId, const std::string& role);
\end{Verbatim}
\textbf{源码：}\path{NDNSF-DistributedInference/cpp/ndnsf-di/NativeProviderSession.hpp}，第 34 行。\par

\Needspace{7\baselineskip}
\textbf{API-070f116bfa58}\quad public\par
\textbf{所属符号：}\code{ndnsf::di::verifyAndUnwrapNativeGrant}\par
\begin{Verbatim}[fontsize=\footnotesize,breaklines,breakanywhere]
NativeGrantVerificationResult
verifyAndUnwrapNativeGrant(const std::string& wireJson,
                           const std::string& authorityPublicKeyRaw,
                           const NativeRecipientKey& recipientKey,
                           const std::string& providerIdentity,
                           const std::string& requestId,
                           std::uint64_t attempt,
                           const std::string& planCoreDigest,
                           const std::string& modelManifestDigest,
                           const std::string& protectionEpoch,
                           std::uint64_t nowMs,
                           const std::string& expectedAuthority = "",
                           const std::string& expectedGrantDigest = "");
\end{Verbatim}
\textbf{源码：}\path{NDNSF-DistributedInference/cpp/ndnsf-di/NativeGrantVerifier.hpp}，第 57 行。\par

完整声明查询：\path{api/target-di-reference.md}。

\Needspace{12\baselineskip}
\section{DI API：生成、KV 与会话延续}
\lead{AC-18 · DI 接口契约}

\subsection{四个 owner 与真正执行入口}
实际生成入口是自由函数 \code{runNativeEpochCoordinator(config)}，不是一个名为 NativeEpochCoordinator 的对象方法。它组织各 role/epoch 执行；\code{QwenGenerationSessionStateMachine} 管理合法动作和终止原因；\code{NativeProviderSession/Runtime} 管理 Provider 解码/会话缓存；\code{NativeConversationCoordinator} 管理跨轮 checkpoint 的本地记录。请求句柄只负责外层等待/取消/观察，不承担这些全部职责。

\subsection{Config 与结果中的关键字段}

\begin{longtable}{@{}p{43mm}p{119mm}@{}}
\toprule
项目 & 行为与调用含义 \\
\midrule\endhead
runtime / plan / assignment / dependencyIo & 构造 config 时提供：runtime 是借用对象，plan/assignment 随 config 保存，I/O 为共享依赖。运行期间这些资源必须有效。 \\
sessionId / requestId / role & 分别识别会话、请求和本地执行角色，localProvider 绑定本地参与者。lineagePlanDigest 来自认证的计划投影。 \\
attemptEpoch / streamEpoch & 默认均为 1；attempt 表示替换尝试，stream 表示流实例，不是每输出一个 token 都增加 streamEpoch。 \\
initialInputs / tokenInputName & 初始张量按名字映射，tokenInputName 默认 input\_ids；\code{stateInputNames/stateOutputNames} 描述模型状态端口。 \\
maxEpochs / sampling & maxEpochs 默认 0，调用方需按有效配置设置界限；samplingMode 默认 Greedy，temperature=0、topK=1、topP=1、repetitionPenalty=1，seed=1750001。不由默认值推断随机采样已验收。 \\
stateIdentityTemplate & 由 adapter/runner 提供不可变状态身份模板，运行时派生前缀/计数相关身份；\code{conversationStateBinding} 指向父轮保留状态，两者不是同一对象。 \\
checkStop / 文本解码 & checkStop 在工作边界报告 Cancelled/Deadline；\code{stableTextDecoder(tokens}, final) 可暂留不稳定后缀，final 时完成 flush。 \\
结果 & finalPayload 为 optional；\code{epochsExecuted/eventsPublished/prefixTokensRecomputed} 是实际计数。finalizedRole 是最后提交的准确角色身份，用于会话状态提升，不应由应用重新拼装。 \\
\bottomrule
\end{longtable}

\subsection{Qwen 状态机：允许的调用}
构造函数先 validate session spec；spec 绑定 \code{candidate/plan/model/artifact}、\code{logicalSessionId/requestId/serviceName}、attempt/token epoch、输入与生成 token 上界、deadline、context/feedback 与角色。下表描述实际方法分支。Preparing 虽在枚举中存在，但 activate 在此实现直接进入 Prefilling，不能据枚举名称补造一条执行链。
\begin{longtable}{@{}p{43mm}p{119mm}@{}}
\toprule
项目 & 行为与调用含义 \\
\midrule\endhead
Created → Selecting & beginSelection，仅 Created 可调用。 \\
Selecting/Rebuilding → Prefilling & activate；其他状态抛 logic\_error。 \\
Prefilling → Decoding & completePrefill；之后才允许 token 计数和停止条件观察。 \\
Decoding → Decoding & \code{completeTokenEpoch(attempt)} 要求 attempt 相同且未达 maxGeneratedTokens，返回增加后的生成计数；旧 attempt 或越界抛 logic\_error。 \\
Decoding → Draining & beginDrain(reason)：Eos/StopSequence 先有对应 observe；MaxTokens 要计数恰好达到上限；ApplicationComplete 无这两项观察要求。 \\
Draining → Completed & complete 必须使用相同 finish reason。claimTerminalResponse 第一次返回 true，后续 false；终态前调用抛异常。它只领取发送权，不负责网络发布。 \\
Decoding → Rebuilding & beginReplacement 增加 attempt，最多到 2；再替换抛异常。生成计数保留，不把已交付前缀清零。 \\
取消 / 期限 / 失败 & cancel 进入 Cancelled，已终态再次调用抛 logic\_error；terminate 要求非终态且 reason 非 None；expireIfDeadlineReached 在尚未到期或已终态时返回 false。 \\
\bottomrule
\end{longtable}

\subsection{两 token 示例与替换}

\begin{Verbatim}[fontsize=\footnotesize,breaklines,breakanywhere]
// 示意：validSpec 已通过模型/角色/摘要与 deadline 校验
QwenGenerationSessionStateMachine state(validSpec);
state.beginSelection();
state.activate();
state.completePrefill();
auto n1 = state.completeTokenEpoch(state.attemptEpoch());
auto n2 = state.completeTokenEpoch(state.attemptEpoch());
// 若生成上限是 2，按 MaxTokens 进入 Draining，再以同一原因 complete
// 终态发布方必须先 claimTerminalResponse()，仅成功领取者发布
\end{Verbatim}

\subsection{token 循环与重复抑制}
epoch coordinator 每角色/轮次调用 NativeModelRunner::run；此路径不使用兼容的单 Provider runStreamed。attempt 2 的 committedPrefixTokenIds 用于重算并核对已提交前缀，抑制重复事件。缓存观测中的 actualNewInputExtent 可能是 token 数，也可能是激活元素数量，不能一律计为 token。checkpointFinalize 是保留会话时的状态补齐阶段，默认 false、上限 32；不发布新的 token 事件。终端文本、事件数量与 KV 包含的前缀必须分别对齐。

\subsection{KV 操作与资源状态}

\begin{longtable}{@{}p{43mm}p{119mm}@{}}
\toprule
项目 & 行为与调用含义 \\
\midrule\endhead
lookupConversationState & 返回 optional<TensorBundle>；空结果不能作为命中。binding 校验错误可能抛异常，不是所有失败都返回空值。 \\
pin / unpin & pin 校验 binding、状态与有效期并增加 pinCount；unpin 递减，计数归零回到 IDLE。pin 不是续期或重新授权。 \\
\code{pauseConversationStateToHost} & 仅匹配、未过期、GPU\_RESIDENT、IDLE 且未 pin 的条目可迁移；Host 空间不足或 adapter 暂停失败返回 false。它改变驻留位置，不新建对话轮次。 \\
\code{prefetchConversationStateToGpu} & 返回 future<bool>，必须等待结果并处理异常；发起预取不表示已在 GPU。\code{cancelConversationStatePrefetch} 与请求取消是不同动作。 \\
\code{releaseConversationState} & 按 binding 释放缓存条目；与 unpin（减少引用保护）不同，也不删除已经保存的会话 journal。 \\
提升缓存 & \code{promoteDecodeStateToConversation} 使用真实执行后的 session/role 和 binding；不可只凭请求号复用另一模型、recipe、前缀或安全域的 KV。 \\
\bottomrule
\end{longtable}

\subsection{从第 N 轮延续到 N+1 轮}
前提是 coordinator 已有父轮记录（或先 restore 加载），Provider 仍有与父 binding 匹配的状态。beginTurn 不能单独启动一个从未登记的新对话：它要求非空 \code{conversationId/serviceName/planRoleMapDigest/parentCheckpointDigest/requestContractDigest}、正 \code{parentContextEpoch/retentionDeadlineMs/request} attempt，并在内存记录中核对父 epoch、摘要、服务和角色映射；不匹配抛 \code{DI_NATIVE_CONVERSATION_PARENT_MISMATCH}。这里仅检查 retention deadline 非零，不自行完成墙钟过期验证。

\subsection{checkpoint 的三个步骤}

\begin{Verbatim}[fontsize=\footnotesize,breaklines,breakanywhere]
// continuation 指向已登记父记录；requestId 为本轮新请求
auto turn = journal.beginTurn(continuation, requestId, 1);
// 真实生成结束后填写 completed：同 request/attempt、tokenIds、
// providerStateDigest，且 complete=true；不能用局部 token 事件冒充
auto checkpoint = journal.prepareCheckpoint(turn, completed);
auto record = journal.commitTurn(turn, checkpoint);
// 下一轮引用 record.checkpoint 的 successor epoch 与摘要
\end{Verbatim}

\subsection{journal 的实际限制}
prepareCheckpoint 拒绝 aborted、request/attempt 不同、未完成、空 token 列表或空 providerStateDigest，报 \code{DI_NATIVE_CONVERSATION_ATTEMPT_INCOMPLETE}。prefixDigest 当前由包含会话/轮次/请求/父摘要/Provider 状态及 token 列表的拼接材料计算，不是仅 token 原始字节摘要。commitTurn 检查父记录未变化，否则报 \code{DI_NATIVE_CONVERSATION_STATE_CHANGED}；写临时文件再 rename 并更新内存。当前没有显式 fsync/输出流失败检查，不能宣称崩溃耐久提交。restore 清空并读取本地 JSON，解析错误可能传播；它不会恢复远端 GPU KV。abortTurn 当前没有持久化或实际修改 turn.aborted 的动作，不能把注释中的 fence 当已实现取消屏障。上述缺口保留为目标改进和 Spec182 待验收边界。

\subsection{完整签名与源码定位}

\Needspace{7\baselineskip}
\textbf{API-d18cce8f82c2}\quad public\par
\textbf{所属符号：}\code{ndnsf::di::runNativeEpochCoordinator}\par
\begin{Verbatim}[fontsize=\footnotesize,breaklines,breakanywhere]
NativeEpochCoordinatorResult
runNativeEpochCoordinator(NativeEpochCoordinatorConfig config);
\end{Verbatim}
\textbf{源码：}\path{NDNSF-DistributedInference/cpp/ndnsf-di/NativeEpochCoordinator.hpp}，第 153 行。\par

\Needspace{7\baselineskip}
\textbf{API-493dc84e5029}\quad public\par
\textbf{所属符号：}\code{ndnsf::di::QwenGenerationSessionStateMachine::completeTokenEpoch}\par
\begin{Verbatim}[fontsize=\footnotesize,breaklines,breakanywhere]
std::uint32_t completeTokenEpoch();
\end{Verbatim}
\textbf{源码：}\path{NDNSF-DistributedInference/cpp/adapters/qwen/QwenGenerationSession.hpp}，第 143 行。\par

\Needspace{7\baselineskip}
\textbf{API-5772f7ae757f}\quad public\par
\textbf{所属符号：}\code{ndnsf::di::QwenGenerationSessionStateMachine::completeTokenEpoch}\par
\begin{Verbatim}[fontsize=\footnotesize,breaklines,breakanywhere]
std::uint32_t completeTokenEpoch(std::uint64_t attemptEpoch);
\end{Verbatim}
\textbf{源码：}\path{NDNSF-DistributedInference/cpp/adapters/qwen/QwenGenerationSession.hpp}，第 144 行。\par

\Needspace{7\baselineskip}
\textbf{API-6fd4334a3745}\quad public\par
\textbf{所属符号：}\code{ndnsf::di::QwenGenerationSessionStateMachine::beginReplacement}\par
\begin{Verbatim}[fontsize=\footnotesize,breaklines,breakanywhere]
void beginReplacement();
\end{Verbatim}
\textbf{源码：}\path{NDNSF-DistributedInference/cpp/adapters/qwen/QwenGenerationSession.hpp}，第 148 行。\par

\Needspace{7\baselineskip}
\textbf{API-30fbe53f3d3e}\quad public\par
\textbf{所属符号：}\code{ndnsf::di::QwenGenerationSessionStateMachine::cancel}\par
\begin{Verbatim}[fontsize=\footnotesize,breaklines,breakanywhere]
void cancel();
\end{Verbatim}
\textbf{源码：}\path{NDNSF-DistributedInference/cpp/adapters/qwen/QwenGenerationSession.hpp}，第 151 行。\par

\Needspace{7\baselineskip}
\textbf{API-855d3ec7afd2}\quad public\par
\textbf{所属符号：}\code{ndnsf::di::QwenGenerationSessionStateMachine::claimTerminalResponse}\par
\begin{Verbatim}[fontsize=\footnotesize,breaklines,breakanywhere]
bool claimTerminalResponse();
\end{Verbatim}
\textbf{源码：}\path{NDNSF-DistributedInference/cpp/adapters/qwen/QwenGenerationSession.hpp}，第 153 行。\par

\Needspace{7\baselineskip}
\textbf{API-1013f1ace7ee}\quad public\par
\textbf{所属符号：}\code{ndnsf::di::NativeProviderSession::lookupConversationState}\par
\begin{Verbatim}[fontsize=\footnotesize,breaklines,breakanywhere]
std::optional<TensorBundle>
  lookupConversationState(const ConversationStateBinding& binding,
                          std::uint64_t nowMs);
\end{Verbatim}
\textbf{源码：}\path{NDNSF-DistributedInference/cpp/ndnsf-di/NativeProviderSession.hpp}，第 51 行。\par

\Needspace{7\baselineskip}
\textbf{API-10751ce7a09d}\quad public\par
\textbf{所属符号：}\code{ndnsf::di::NativeProviderSession::pauseConversationStateToHost}\par
\begin{Verbatim}[fontsize=\footnotesize,breaklines,breakanywhere]
bool
  pauseConversationStateToHost(const ConversationStateBinding& binding,
                               std::uint64_t nowMs);
\end{Verbatim}
\textbf{源码：}\path{NDNSF-DistributedInference/cpp/ndnsf-di/NativeProviderSession.hpp}，第 55 行。\par

\Needspace{7\baselineskip}
\textbf{API-6b3bff150575}\quad public\par
\textbf{所属符号：}\code{ndnsf::di::NativeProviderSession::prefetchConversationStateToGpu}\par
\begin{Verbatim}[fontsize=\footnotesize,breaklines,breakanywhere]
std::future<bool>
  prefetchConversationStateToGpu(const ConversationStateBinding& binding,
                                  std::uint64_t nowMs);
\end{Verbatim}
\textbf{源码：}\path{NDNSF-DistributedInference/cpp/ndnsf-di/NativeProviderSession.hpp}，第 59 行。\par

\Needspace{7\baselineskip}
\textbf{API-daca2274c659}\quad public\par
\textbf{所属符号：}\code{ndnsf::di::NativeProviderSession::pinConversationState}\par
\begin{Verbatim}[fontsize=\footnotesize,breaklines,breakanywhere]
bool
  pinConversationState(const ConversationStateBinding& binding,
                       std::uint64_t nowMs);
\end{Verbatim}
\textbf{源码：}\path{NDNSF-DistributedInference/cpp/ndnsf-di/NativeProviderSession.hpp}，第 66 行。\par

\Needspace{7\baselineskip}
\textbf{API-32c8162259ef}\quad public\par
\textbf{所属符号：}\code{ndnsf::di::NativeProviderSession::releaseConversationState}\par
\begin{Verbatim}[fontsize=\footnotesize,breaklines,breakanywhere]
bool
  releaseConversationState(const ConversationStateBinding& binding);
\end{Verbatim}
\textbf{源码：}\path{NDNSF-DistributedInference/cpp/ndnsf-di/NativeProviderSession.hpp}，第 73 行。\par

\Needspace{7\baselineskip}
\textbf{API-1e7281b52fd5}\quad public\par
\textbf{所属符号：}\code{ndnsf::di::NativeConversationCoordinator::beginTurn}\par
\begin{Verbatim}[fontsize=\footnotesize,breaklines,breakanywhere]
NativeConversationTurn beginTurn(const NativeConversationContinuation& continuation,
                                   std::string requestId,
                                   std::uint64_t attempt = 1) const;
\end{Verbatim}
\textbf{源码：}\path{NDNSF-DistributedInference/cpp/ndnsf-di/NativeConversationCoordinator.hpp}，第 73 行。\par

\Needspace{7\baselineskip}
\textbf{API-a550281a7724}\quad public\par
\textbf{所属符号：}\code{ndnsf::di::NativeConversationCoordinator::prepareCheckpoint}\par
\begin{Verbatim}[fontsize=\footnotesize,breaklines,breakanywhere]
NativeConversationCheckpoint prepareCheckpoint(
    const NativeConversationTurn& turn, const NativeCompletedAttempt& completed) const;
\end{Verbatim}
\textbf{源码：}\path{NDNSF-DistributedInference/cpp/ndnsf-di/NativeConversationCoordinator.hpp}，第 77 行。\par

\Needspace{7\baselineskip}
\textbf{API-8c7c62356113}\quad public\par
\textbf{所属符号：}\code{ndnsf::di::NativeConversationCoordinator::commitTurn}\par
\begin{Verbatim}[fontsize=\footnotesize,breaklines,breakanywhere]
NativeConversationRecord commitTurn(const NativeConversationTurn& turn,
                                      const NativeConversationCheckpoint& checkpoint);
\end{Verbatim}
\textbf{源码：}\path{NDNSF-DistributedInference/cpp/ndnsf-di/NativeConversationCoordinator.hpp}，第 79 行。\par

\Needspace{7\baselineskip}
\textbf{API-64cbde6b1261}\quad public\par
\textbf{所属符号：}\code{ndnsf::di::NativeConversationCoordinator::restore}\par
\begin{Verbatim}[fontsize=\footnotesize,breaklines,breakanywhere]
void restore(const std::filesystem::path& journalRoot);
\end{Verbatim}
\textbf{源码：}\path{NDNSF-DistributedInference/cpp/ndnsf-di/NativeConversationCoordinator.hpp}，第 81 行。\par

完整声明查询：\path{api/target-di-reference.md}。

\Needspace{12\baselineskip}
\section{UAV API：飞控边界与操作权限}
\lead{AC-19 · UAV 接口契约}

\subsection{内部接口与线载荷}
FlightControllerBackend 是 Drone 应用内部接口，sendMavlink(frame, commandName) 接收原始 MAVLink 字节与命令标签，返回 Fields 键值结果。Fields 不提供统一静态类型/单位检查；调用者必须遵循 UavProtocol 与当前服务 handler 的命令字段解析，不能把任意 map 当已经过认证的命令。AC-19 所列为接口声明，Mock/UDP/串口行为要按实际 backend 分开。

\subsection{完整命令路径}
GS 从 UI 状态和操作选择构造命令，经 /UAV/MAVLink/Execute 调用目标 Drone；Core 检查服务权限、token 和认证绑定，UAV 路径处理操作 lease，应用/飞控处理 readiness 与执行。收到网络 Response、backend 发送字节和飞控 COMMAND\_ACK 是三个事件；UI 应保留失败所处阶段，不能把一次本地 send 成功显示成车辆已经执行。

\subsection{三类状态的例子}

\begin{longtable}{@{}p{43mm}p{119mm}@{}}
\toprule
项目 & 行为与调用含义 \\
\midrule\endhead
Core 授权失败 & 尚未获得服务执行资格；修复身份/策略，不能重新解释为飞控忙。 \\
OperatorAuthorityLease 失效 & 操作员对目标/范围/期限的应用授权不成立；issuer 不是 ServiceController，lease 版本不是 Controller epoch。 \\
readiness 不满足 & 例如应用所需新鲜遥测缺失；lease 有效不意味着所有飞行动作都可执行。 \\
Mock 与设备 backend & Mock 返回可用于界面/状态机检查的结果；UDP/串口涉及线程、链路、ACK 和设备错误，其成功不能从 Mock 外推。 \\
\bottomrule
\end{longtable}

\subsection{扩展边界}
新 backend 实现能力、遥测、命令结果与停止行为，保留实际错误。GS Qt 回调不能延长已销毁设备对象的生命周期。当前仍缺一份覆盖所有命令的规范字段/ACK 表，不能把本章当完整飞控线协议规范；TG-05 规定类型化与兼容转换的收敛方向。

\subsection{当前 Execute 字段与 backend 返回}
Fields 的实际类型是 map<string,string>。Drone Execute handler 读取 command（缺省 unknown）和 mavlink\_hex（缺省空字符串），hexDecode 后直接调用 backend->sendMavlink；返回附加 backend 和 drone\_id，以 accepted 字符串是否为 true 决定外层响应成功。此 handler 内未见 OperatorAuthorityLease 或 readiness 的独立重验；GS 的操作门与 UI 拒绝不能作为该网络入口已强制执行这两项检查的证据。
\begin{longtable}{@{}p{43mm}p{119mm}@{}}
\toprule
项目 & 行为与调用含义 \\
\midrule\endhead
accepted / ack\_result & Mock 返回 \code{mock-accepted/mock-rejected}。设备 backend 连接/发送失败分别返回 \code{connect-failed/send-failed}；其他命令以 accepted 或 in-progress ACK 判为 accepted=true。 \\
manual\_control & 成功发送后仅收集约 5 ms 遥测，设置 \code{manual-control-forwarded}，并据转发判为接受，不等同等待 COMMAND\_ACK。 \\
其他设备命令 & sendMavlink 传入 700 ms ACK 等待预算；整次调用还含连接和发送，不保证总耗时严格等于 700 ms。内部持 socket mutex，可能阻塞调用线程。 \\
forwarded\_bytes & 记录帧大小；send-failed 分支也填 frame.size()，不是实际成功发送字节的可靠统计。 \\
ack\_source / fc\_state & ack\_source 是 mock 或实际 transport；fc\_state 及附加遥测由后端产生，字段集合并非所有 backend 完全一致。 \\
\bottomrule
\end{longtable}

\subsection{完整签名与源码定位}

\Needspace{7\baselineskip}
\textbf{API-b1e46a3def66}\quad public\par
\textbf{所属符号：}\code{FlightControllerBackend::sendMavlink}\par
\begin{Verbatim}[fontsize=\footnotesize,breaklines,breakanywhere]
virtual Fields sendMavlink(const std::vector<uint8_t>& frame,
                             const std::string& commandName) = 0;
\end{Verbatim}
\textbf{源码：}\path{NDNSF-UAV-APP/drone/DroneServiceContainer.inc.hpp}，第 8 行。\par

完整声明查询：\path{api/target-uav-reference.md}。

\Needspace{12\baselineskip}
\section{UAV API：任务、补偿与终端报告}
\lead{AC-20 · UAV 接口契约}

\subsection{任务和调用各自存活多久}
UavMissionSession 持有长生命周期 mission；part 是可独立完成的工作单元，incident/job 是有限分析调用。assignPart、markPartExecuting、completePart、markPartMissing 维护 part。Core requestId/token 属于每次 invocation，新 attempt 不沿用旧调用令牌。

\subsection{两个 part 的补偿例子}
任务含 P0、P1：P0 已完成并有有效结果，P1 失联进入 missing。compensateMissingParts 只为 P1 形成补偿工作，不重跑 P0。补偿使用新的 attempt/调用绑定；若旧 P1 结果晚到，acceptTerminalReport 仍要匹配当前任务、job、attempt、计划与 terminal owner，不能覆盖新 attempt。全部必要 part 有有效结果才是 mission 完成。

\subsection{返回、恢复与资源}
bool 表示本地迁移/报告接受，reason 指针可接收诊断；false 不表示车辆已停止。snapshot/restore 保存/载入业务状态，\code{reconcileVehicleAndStreams} 再与实际车辆/流对照；不能恢复一个 Active 标记就继续发命令。incident 的 \code{retryAnalysis/retriesUsed} 另有有界重试，不由 mission 无限重试代替。

\subsection{完整签名与源码定位}

\Needspace{7\baselineskip}
\textbf{API-9439cf6844b2}\quad public\par
\textbf{所属符号：}\code{ndnsf::examples::uav::UavMissionSession::compensateMissingParts}\par
\begin{Verbatim}[fontsize=\footnotesize,breaklines,breakanywhere]
bool compensateMissingParts(std::string* reason = nullptr);
\end{Verbatim}
\textbf{源码：}\path{NDNSF-UAV-APP/shared/UavMissionSession.hpp}，第 49 行。\par

\Needspace{7\baselineskip}
\textbf{API-8bc23c85a02b}\quad public\par
\textbf{所属符号：}\code{ndnsf::examples::uav::UavMissionSession::acceptTerminalReport}\par
\begin{Verbatim}[fontsize=\footnotesize,breaklines,breakanywhere]
bool acceptTerminalReport(const UavTerminalReport& report,
                            std::string* reason = nullptr);
\end{Verbatim}
\textbf{源码：}\path{NDNSF-UAV-APP/shared/UavMissionSession.hpp}，第 60 行。\par

\Needspace{7\baselineskip}
\textbf{API-f10389b9d21d}\quad public\par
\textbf{所属符号：}\code{ndnsf::examples::uav::UavIncidentCoordinator::acceptReport}\par
\begin{Verbatim}[fontsize=\footnotesize,breaklines,breakanywhere]
bool acceptReport(const UavTerminalReport& report, uint64_t nowMs,
                    std::string* reason = nullptr);
\end{Verbatim}
\textbf{源码：}\path{NDNSF-UAV-APP/ground-station/UavIncidentCoordinator.hpp}，第 62 行。\par

完整声明查询：\path{api/target-uav-reference.md}。

\Needspace{12\baselineskip}
\section{UAV API：视频管线、遥测与录像}
\lead{AC-21 · UAV 接口契约}

\subsection{队列输入与准确返回}
BoundedLatestFrameQueue 构造要求 maxFrames 和 maxBytes 都为正，否则 invalid\_argument。push 接收含字节的 UavVideoFrame；单帧超过字节容量返回 false 并记录 \code{frame-exceeds-byte-capacity}；否则先丢最旧帧直到数量和字节均容纳新帧，再入队返回 true。因此 true 也可能伴随历史帧被丢弃。

\subsection{popLatest 的例子}
队列依次存有帧 10、11、12，popLatest 只返回 12，把 10、11 计为 dropped，并清空队列和字节计数；空队列返回 nullopt。此策略适合及时显示，不能用于要求逐段完整的制品下载或完整录像。push/pop/clear 的内部 mutex 保护队列，不自动保护外部解码器。

\subsection{遥测值与时间}

\begin{longtable}{@{}p{43mm}p{119mm}@{}}
\toprule
项目 & 行为与调用含义 \\
\midrule\endhead
地理/速度/电量 & latitudeE7/longitudeE7 为角度的 10\textasciicircum{}7 缩放，altitudeMm 为毫米，groundSpeedMmps 为毫米每秒，batteryPermille 为千分比。 \\
sampleId / sourceTimestampNs & 分别用于样本身份和源时间；receivedTimestampNs 是接收时间，不能与源时间混作同一观测。 \\
admit 返回 & \code{valid/stateAdvanced/newSample/duplicate/outOfOrder/ageNs/reason} 分开解释；格式有效的旧样本不一定推进最新状态。 \\
声学块 & \code{CompleteAcousticBlockAdmission} 要求领域完整性和截止期；Core 的 FEC/签名恢复不直接等价于声学业务可用。 \\
\bottomrule
\end{longtable}

\subsection{捕获到录像的分工}
capture/编码产生带 session、frameId、PTS、keyframe 等元数据的帧，Core 流负责命名传输，GS 解码后经最新帧队列显示。录像路径保存原签名/保护 packet 与应用 manifest 到 Repo，记录 gap、retention 与游标；直播显示丢帧和录像数据缺口分别统计。相机链已有具体入口，不能据 DataProductReference 推定所有遥测/检测结果已自动录制。

\subsection{完整签名与源码定位}

\Needspace{7\baselineskip}
\textbf{API-a064d74ab888}\quad public\par
\textbf{所属符号：}\code{ndnsf::examples::uav::BoundedLatestFrameQueue::push}\par
\begin{Verbatim}[fontsize=\footnotesize,breaklines,breakanywhere]
bool push(UavVideoFrame frame);
\end{Verbatim}
\textbf{源码：}\path{NDNSF-UAV-APP/shared/UavVideoPipeline.hpp}，第 65 行。\par

\Needspace{7\baselineskip}
\textbf{API-11ef609f2477}\quad public\par
\textbf{所属符号：}\code{ndnsf::examples::uav::BoundedLatestFrameQueue::popLatest}\par
\begin{Verbatim}[fontsize=\footnotesize,breaklines,breakanywhere]
std::optional<UavVideoFrame> popLatest();
\end{Verbatim}
\textbf{源码：}\path{NDNSF-UAV-APP/shared/UavVideoPipeline.hpp}，第 66 行。\par

\Needspace{7\baselineskip}
\textbf{API-730501b40b2c}\quad public\par
\textbf{所属符号：}\code{ndnsf::examples::uav::LatestTelemetryAdmission::admit}\par
\begin{Verbatim}[fontsize=\footnotesize,breaklines,breakanywhere]
TelemetryAdmissionResult admit(const std::vector<uint8_t>& wire,
                                 uint64_t receivedTimestampNs);
\end{Verbatim}
\textbf{源码：}\path{NDNSF-UAV-APP/shared/UavSensorStreams.hpp}，第 57 行。\par

完整声明查询：\path{api/target-uav-reference.md}。

\Needspace{12\baselineskip}
\section{UAV API：检测与多视角识别扩展}
\lead{AC-22 · UAV 接口契约}

\subsection{输入对象与责任}
MultiViewRecognitionJob 绑定任务/attempt、目标与精确视图引用；model profile 指定模型输入要求，verified views 提供已验证视图，terminalOwner 指定结果归属。executeMultiViewJob 返回 \code{MultiViewExecutionResult}；入口先通过 \code{validateMultiViewJobAndViews} 校验任务、视角数量/来源、时间窗口和模型绑定，不能让算法自行忽略缺视角。

\subsection{两个视角的例子}
Drone A 和 B 在同一事件窗口提供两份已验证画面引用，job 指向这两份精确对象。选定算法消费对应视图并产生结果，\code{UavIncidentCoordinator::acceptMultiViewResult} 再校验当前 job/attempt/owner。把 A 的同一画面重复两次不应作为两架独立 Drone 的证据；单视角基线另行构造输入。

\subsection{算法扩展与错误}
\code{IMultiViewRecognitionAlgorithm} 是 C++ 扩展点；默认 \code{DeterministicMultiViewAlgorithm} 是确定性研究/测试夹具。真实 ONNX CPUExecutionProvider 路径在 \code{tools/mvcnn_onnx_worker.py}，必须显式核对调用方接线。缺视角、模型/摘要不符、任务过期、错误 owner 与模型执行失败保留各自原因。GS 的单帧 YOLO、UAV 多视角 job、NDNSF-DI 规划是三个入口，不能混称一次分布式推理。

\subsection{完整签名与源码定位}

\Needspace{7\baselineskip}
\textbf{API-1da9ec7348c6}\quad public\par
\textbf{所属符号：}\code{ndnsf::examples::uav::executeMultiViewJob}\par
\begin{Verbatim}[fontsize=\footnotesize,breaklines,breakanywhere]
MultiViewExecutionResult
executeMultiViewJob(const MultiViewRecognitionJob& job,
                    const MultiViewModelProfile& profile,
                    const std::vector<VerifiedMultiViewInput>& views,
                    const ndn::Name& terminalOwner,
                    const IMultiViewRecognitionAlgorithm& algorithm =
                      DeterministicMultiViewAlgorithm{});
\end{Verbatim}
\textbf{源码：}\path{NDNSF-UAV-APP/shared/UavMultiViewRecognition.hpp}，第 195 行。\par

\Needspace{7\baselineskip}
\textbf{API-01be92408d26}\quad public\par
\textbf{所属符号：}\code{ndnsf::examples::uav::UavIncidentCoordinator::acceptMultiViewResult}\par
\begin{Verbatim}[fontsize=\footnotesize,breaklines,breakanywhere]
bool acceptMultiViewResult(const FusedRecognitionResult& result, uint64_t nowMs,
                             std::string* reason = nullptr);
\end{Verbatim}
\textbf{源码：}\path{NDNSF-UAV-APP/ground-station/UavIncidentCoordinator.hpp}，第 66 行。\par

完整声明查询：\path{api/target-uav-reference.md}。

\Needspace{12\baselineskip}
\section{API 扩展与跨模块调用检查}
\lead{AC-23 · Core 接口契约}

\subsection{新增服务的最小闭环}
以 echo 服务为例：应用定义可序列化请求/响应，Provider 在 ACK 阶段只判断资格，Selection 通过后执行并返回；保存 scoped registration；User 用 AC-03 普通调用保存 requestId 并处理响应/超时；停服时关闭注册并排空工作。跨网络测试应覆盖无权限、无候选、错误响应与迟到回调，而不仅是本地 handler 直接调用。Core 不为这个服务增加领域字段。

\subsection{新增 DI 策略的最小闭环}
实现 AC-15 的 identity/enumerate 或 proposeRoles，用固定模型图和认证报价生成候选；交由共享 validator，准备/发布/再绑定后封存；将实现登记到实际使用的 registry/配置；验证默认调用链确实选中该策略。只增加类、只通过 mock 端口测试或只生成更多 API ID，都不能表示该功能接入产品。

\subsection{设计维护的交付单元}
API 变化同时登记输入/输出字段、单位、错误、状态、线程与资源、源码位置和验证证据。按 MANAGEMENT.md 更新当前设计；目标独立维护，PLANNED 不写成 existing。生成器验证声明，人工审阅验证解释能否指导调用；每个 Spec 在 spec-design-changes.md 记录前后行为，不用文件数量或模板段落宣称完整。

完整声明查询：\path{api/target-core-reference.md}。
